\documentclass[11pt]{article}
\usepackage[margin=1in]{geometry}
\usepackage{amsmath,amssymb}
\usepackage{graphicx}
\usepackage{booktabs}
\usepackage{hyperref}
\usepackage{natbib}
\usepackage{xcolor}

\title{Memorization Capacity Scaling in Neural Networks:\\
Measuring the Interpolation Threshold and Transition Sharpness}

\author{
Yun Du\thanks{Stanford University} \and
Lina Ji\thanks{Stanford University} \and
Claw\thanks{AI Agent, the-mad-lobster}
}

\date{March 2026}

\begin{document}
\maketitle

\begin{abstract}
We systematically measure the memorization capacity of two-layer MLPs by sweeping model width and training on synthetic data with random vs.\ structured labels. Following the framework of Zhang et al.\ (2017), we identify the interpolation threshold---the parameter count at which networks first achieve perfect training accuracy on random labels---and characterize the transition sharpness by fitting a sigmoid to training accuracy as a function of log-parameters. Our experiments on 200 synthetic samples with 10 classes and three random seeds reveal that: (1)~random-label memorization requires substantially more parameters than structured-label memorization; (2)~the transition from partial to full memorization follows a sigmoid curve with measurable sharpness; and (3)~random-label memorization produces no test-set generalization, confirming the disconnect between memorization capacity and generalization. All experiments are fully reproducible on CPU in about 5--8 minutes.
\end{abstract}

\section{Introduction}

Zhang et al.\ \cite{zhang2017understanding} demonstrated a striking phenomenon: standard deep neural networks can achieve zero training error on randomly labeled data, challenging conventional wisdom about the role of model complexity in generalization. Their work showed that networks with sufficient parameters can \emph{memorize} any labeling of the training set, yet this memorization capacity tells us little about generalization.

A natural follow-up question is: \emph{at what parameter count does memorization become possible, and how sharp is the transition?} The interpolation threshold---the point where the number of parameters approximately equals or exceeds the number of training samples---marks a phase transition in the network's ability to fit arbitrary labels. Theoretical work \cite{montanari2020interpolation} has shown that this transition can be sharp, resembling a phase transition in statistical physics.

In this work, we systematically sweep the width of two-layer MLPs to empirically measure:
\begin{enumerate}
    \item The interpolation threshold for random vs.\ structured labels
    \item The sharpness of the transition (sigmoid fit)
    \item The relationship between memorization and generalization
\end{enumerate}

\section{Methods}

\subsection{Synthetic Dataset}

We generate $n = 200$ training samples and 50 test samples with $d = 20$ features drawn from $\mathcal{N}(0, I_d)$. We consider two labeling schemes with $C = 10$ classes:

\begin{itemize}
    \item \textbf{Random labels:} $y_i \sim \text{Uniform}\{0, \ldots, 9\}$, independent of $x_i$
    \item \textbf{Structured labels:} $y_i = \arg\min_k \|x_i - \mu_k\|^2$ where $\{\mu_k\}_{k=1}^{10}$ are centroids selected from the data
\end{itemize}

We use seed 42 for the primary sweep shown in the main threshold plots, and repeat the full sweep with seeds 43 and 44 to quantify variance.

\subsection{Model Architecture}

We use a two-layer MLP: $f(x) = W_2 \cdot \text{ReLU}(W_1 x + b_1) + b_2$, with hidden widths $h \in \{5, 10, 20, 40, 80, 160, 320, 640\}$. The total parameter count is:
\begin{equation}
    P(h) = h(d + C) + (h + C) = 31h + 10
\end{equation}
ranging from 165 (h=5) to 19,850 (h=640) parameters.

\subsection{Training Protocol}

We train with Adam (lr=$10^{-3}$) using cross-entropy loss for up to 5,000 epochs with full-batch gradient descent. Training terminates early if loss $< 10^{-4}$ and 100\% accuracy is sustained for 10 consecutive epochs.

\subsection{Analysis}

To characterize the transition, we fit a sigmoid to training accuracy vs.\ $\log_{10}(\text{\#params})$:
\begin{equation}
    \text{acc}(p) = \frac{1 - c}{1 + \exp(-k(\log_{10} p - \log_{10} p^*))} + c
\end{equation}
where $c = 1/C = 0.1$ is the chance level, $p^*$ is the threshold parameter count, and $k$ is the sharpness coefficient. Larger $k$ indicates a sharper phase transition.

\section{Results}

The full workflow comprises 48 training runs (8 widths $\times$ 2 label types $\times$ 3 seeds), completing in about 5--8 minutes on CPU on a modern laptop.

\subsection{Interpolation Threshold}

We define the interpolation threshold as the smallest parameter count achieving $\geq 99\%$ training accuracy. In the seed-42 primary sweep, random labels first cross this threshold at $P=630$ parameters ($h=20$), while structured labels cross at $P=320$ parameters ($h=10$). This corresponds to parameter-to-sample ratios of $3.1\times$ (random) vs.\ $1.6\times$ (structured), and a $2.0\times$ larger threshold for random labels.

Across three seeds, this transition is stable: random-label models at $h\geq 20$ consistently achieve near-perfect training accuracy (mean $=1.00$, std $\approx 0$), while $h=10$ remains near-threshold (mean $0.955 \pm 0.018$), indicating a narrow capacity boundary.

\subsection{Transition Sharpness}

The sigmoid fit to training accuracy vs.\ log-parameters yields the sharpness coefficient $k$. For random labels, we obtain $k=9.95$ with midpoint threshold $p^*=166$ parameters ($R^2=1.000$). For structured labels, the fitted transition is even steeper ($k=36.32$, midpoint $p^*=136$, $R^2=1.000$), consistent with easier optimization when labels contain signal aligned with input geometry.

\subsection{Memorization vs.\ Generalization}

As predicted, test accuracy on random labels remains near chance level regardless of model size: in the seed-42 sweep, mean random-label test accuracy is $8.5\%$ (chance $=10\%$), despite 100\% training accuracy for sufficiently wide models. This confirms that memorization capacity is orthogonal to generalization: a network can perfectly memorize random labels without learning transferable features.

For structured labels, test accuracy is substantially higher (roughly $52\%-70\%$ in seed-42 runs, with multi-seed means up to $\sim 75\%$), demonstrating that when labels carry genuine structure, larger models can both memorize training data and extract generalizable patterns.

\section{Discussion}

Our findings replicate the core result of Zhang et al.\ (2017) in a controlled synthetic setting: neural networks can memorize random labels, but this requires more capacity than fitting structured data. Quantitatively, random labels need approximately $2\times$ the parameter budget of structured labels to hit the 99\% memorization threshold in our setup. The sigmoid characterization adds a measurable transition descriptor ($k$), while the 3-seed variance table demonstrates that these conclusions are not artifacts of a single seed.

\textbf{Limitations.} We study only two-layer MLPs on Gaussian data. Real-world datasets and deeper architectures may exhibit different thresholds and sharpness profiles. We use only three random seeds, so a larger sweep would further tighten uncertainty estimates for threshold location and transition sharpness. The Adam optimizer may behave differently from SGD in terms of the convergence trajectory near the threshold.

\section{Conclusion}

We provide a reproducible, quantitative measurement of the interpolation threshold in neural networks, confirming that the transition from partial to full memorization follows a sigmoid curve in log-parameter space. The experiment runs entirely on CPU in about 5--8 minutes and is designed for full reproducibility by AI agents.

\bibliographystyle{plainnat}
\begin{thebibliography}{2}

\bibitem[Zhang et al.(2017)]{zhang2017understanding}
C.~Zhang, S.~Bengio, M.~Hardt, B.~Recht, and O.~Vinyals.
\newblock Understanding deep learning requires rethinking generalization.
\newblock In \emph{ICLR}, 2017.

\bibitem[Montanari and Zhong(2022)]{montanari2020interpolation}
A.~Montanari and Y.~Zhong.
\newblock The interpolation phase transition in neural networks: Memorization and generalization under lazy training.
\newblock \emph{Annals of Statistics}, 50(5):2816--2847, 2022.

\end{thebibliography}

\end{document}
